\section{Interpretación integrada}
\label{sec:interpretacion-integrada}

La pregunta central del trabajo fue:

\begin{quote}
    \itshape
    ¿Qué factores explican la probabilidad de que un solicitante sea
    clasificado como mal riesgo crediticio?
\end{quote}

La evidencia obtenida muestra que el riesgo crediticio no se concentra
en una única característica. La clasificación depende de una combinación
entre las condiciones del préstamo, la liquidez disponible, los
antecedentes financieros y algunos indicadores de estabilidad económica.

El término ``explican'' se utiliza en un sentido estadístico y
predictivo. Los resultados permiten identificar variables asociadas con
cambios en la probabilidad estimada de mal crédito, pero no demuestran
que esas variables causen directamente el resultado observado.

\subsection{Respuesta a la pregunta central}

Dentro de las variables numéricas, la duración del préstamo constituye
el factor con la asociación más consistente. Los créditos clasificados
como malos presentan plazos mayores en el análisis descriptivo, una
diferencia estandarizada cercana a moderada en la comparación
frecuentista y un coeficiente positivo en la regresión logística.

En el modelo multivariable, cada mes adicional de duración se asocia con
un incremento aproximado del \(3.9\,\%\) en las \textit{odds} de mal
crédito, manteniendo constantes los demás predictores. El análisis
bayesiano confirmó este resultado:

\[
\operatorname{E}
\left[
    \beta_{\mathrm{dur}}
    \mid
    \mathrm{datos}
\right]
=
0.04084,
\]

\[
\operatorname{ICr}_{95\,\%}
=
[0.01902,\;0.06275],
\]

y:

\[
\Pr
\left(
    \beta_{\mathrm{dur}}>0
    \mid
    \mathrm{datos}
\right)
\approx
1.
\]

Por tanto, la duración es el principal factor numérico asociado con un
mayor riesgo dentro de la especificación utilizada.

La carga de la cuota también mantiene una asociación condicional
relevante. Una unidad adicional en la escala ordinal de
\variable{tasa_cuota} se relaciona con un aumento aproximado del
\(22.8\,\%\) en las \textit{odds} de mal crédito:

\[
\mathrm{OR}=1.2283.
\]

Este resultado es coherente con la interpretación financiera de que una
mayor proporción de recursos comprometida en el pago puede reducir el
margen disponible frente a gastos inesperados. No obstante, la variable
representa niveles ordinales y no porcentajes exactos del ingreso.

Entre las variables categóricas, el estado de la cuenta corriente, el
nivel de ahorros, algunas categorías del historial crediticio, la
antigüedad laboral y el tipo de vivienda aportan información adicional.
En términos generales, los perfiles con mejores señales de liquidez o
estabilidad presentan menores \textit{odds} de mal crédito respecto de
sus categorías de referencia.

La interpretación de estas categorías requiere cautela. Algunos códigos
históricos combinan situaciones diferentes y ciertas asociaciones pueden
parecer contraintuitivas debido a la categoría utilizada como referencia.
En consecuencia, no deben convertirse directamente en reglas generales
de aprobación o rechazo.

\subsection{Diferencia entre asociación marginal y condicional}

Uno de los resultados más importantes del análisis es la diferencia entre
una asociación observada por separado y una asociación que permanece al
controlar por otras variables.

El monto del crédito mostró inicialmente:

\begin{itemize}
    \item valores medios y medianos superiores entre los malos créditos;

    \item una diferencia estadísticamente significativa mediante
    Mann--Whitney;

    \item una asociación positiva, aunque pequeña, con la clase de
    riesgo.
\end{itemize}

Sin embargo, dentro de la regresión logística:

\[
\mathrm{OR}_{\log(\mathrm{monto})}
=
1.0658,
\]

con:

\[
p=0.7238.
\]

El intervalo del \(95\,\%\) contiene el valor uno. Por tanto, no existe
evidencia suficiente de una asociación independiente del monto después
de considerar duración, estado de cuenta, historial, ahorros, carga de
la cuota y vivienda.

Este cambio es coherente con la relación observada entre monto y
duración:

\[
\rho_{\mathrm{Spearman}}
\approx
0.62.
\]

Los créditos de mayor monto tienden a concederse por plazos más largos.
En consecuencia, parte de la asociación marginal del monto con el riesgo
parece estar capturada por la duración y por otras condiciones
financieras.

La edad presenta un comportamiento parecido. Los malos créditos
corresponden descriptivamente a solicitantes algo más jóvenes, pero la
edad no conserva evidencia estadística dentro del modelo logístico:

\[
\mathrm{OR}_{\mathrm{edad}}
=
0.9972,
\qquad
p=0.7798.
\]

Por tanto, ni el monto ni la edad deben considerarse factores
condicionales principales en la especificación final.

\subsection{Propósito del préstamo como factor de segmentación}

La prueba chi-cuadrado identificó una asociación entre el propósito y la
clase de riesgo:

\[
\chi^2(9)
=
33.356,
\qquad
p=0.00012.
\]

La magnitud estimada mediante V de Cramér fue:

\[
V=0.183.
\]

Esto indica una asociación global estadísticamente significativa, pero
pequeña. Algunos propósitos presentan tasas observadas superiores al
promedio, mientras que otros muestran tasas menores.

Sin embargo, el propósito debe entenderse principalmente como una
variable de segmentación. Las tasas de algunas categorías se basan en
pocos registros y pueden ser inestables. Además, la prueba global no
controla por duración, monto, liquidez ni situación financiera.

Por ello, el propósito aporta contexto sobre el tipo de operación, pero
no constituye por sí solo una base suficiente para determinar el riesgo
individual de un solicitante.

\subsection{Síntesis de la evidencia por factor}

La Tabla~\ref{tab:evidencia-integrada-factores} resume la relación entre
los principales factores y el riesgo.

\begin{table}[H]
    \centering
    \caption{Síntesis integrada de los factores asociados con el riesgo
    crediticio.}
    \label{tab:evidencia-integrada-factores}

    \small
    \begin{tabularx}{\textwidth}{
        @{}
        >{\raggedright\arraybackslash}p{2.7cm}
        >{\raggedright\arraybackslash}p{4.1cm}
        >{\raggedright\arraybackslash}p{4.0cm}
        X
        @{}
    }
        \toprule
        Factor
        & Evidencia marginal
        & Evidencia condicional
        & Interpretación integrada \\
        \midrule

        Duración
        & Los malos créditos tienen plazos mayores; \(d=0.480\).
        & \(\mathrm{OR}=1.0390\) por mes; posterior completamente
          concentrada en valores positivos.
        & Factor numérico más consistente del análisis. \\

        Tasa de la cuota
        & Mayor presencia de niveles altos entre los malos créditos.
        & \(\mathrm{OR}=1.2283\), con intervalo por encima de uno.
        & Una mayor carga financiera se asocia con mayor riesgo. \\

        Monto
        & Diferencia significativa entre clases; efecto pequeño.
        & \(\mathrm{OR}=1.0658\), \(p=0.7238\).
        & Su asociación marginal no permanece después del control
          multivariable. \\

        Estado de cuenta
        & Diferencias visibles entre perfiles de liquidez.
        & Varias categorías presentan \(\mathrm{OR}<1\) frente al saldo
          negativo.
        & La liquidez corriente aporta una señal condicional importante. \\

        Ahorros
        & Los perfiles con mayor ahorro muestran menor riesgo observado.
        & \texttt{A64}: \(\mathrm{OR}=0.1617\) frente a menos de
          \(100\) DM.
        & Una mayor capacidad de ahorro se asocia con menor riesgo. \\

        Historial crediticio
        & Las tasas cambian entre categorías.
        & Algunas categorías presentan asociaciones significativas
          respecto de \texttt{A30}.
        & Aporta información, pero depende fuertemente de la referencia
          y de la codificación histórica. \\

        Vivienda
        & La distribución de las clases varía según el régimen de
          vivienda.
        & Vivienda propia:
          \(\mathrm{OR}=0.5868\) frente al alquiler.
        & La estabilidad residencial se asocia con menores
          \textit{odds}. \\

        Propósito
        & Asociación global significativa, \(V=0.183\).
        & No se utilizó como predictor en la especificación logística
          final.
        & Útil para segmentación, con capacidad individual limitada. \\

        Edad
        & Los malos créditos son ligeramente más jóvenes.
        & \(\mathrm{OR}=0.9972\), \(p=0.7798\).
        & La diferencia marginal no persiste en el modelo conjunto. \\

        \bottomrule
    \end{tabularx}
\end{table}

\subsection{Respuesta a las preguntas secundarias}

Las preguntas específicas del trabajo pueden responderse de manera
integrada.

\subsubsection{¿Cuál es la proporción global de malos créditos?}

La proporción observada fue:

\[
\widehat{p}
=
0.300.
\]

El intervalo de Wilson del \(95\,\%\) fue:

\[
[0.2724,\;0.3291],
\]

mientras que el intervalo bootstrap fue:

\[
[0.2720,\;0.3290].
\]

El modelo Beta--Binomial produjo intervalos creíbles prácticamente
idénticos bajo los tres priors evaluados. La tasa se estima, por tanto,
cerca del \(30\,\%\), con una incertidumbre aproximada de tres puntos
porcentuales a cada lado.

\subsubsection{¿La tasa difiere según el propósito?}

Sí. La prueba chi-cuadrado rechazó la hipótesis de independencia entre
propósito y clase de riesgo. Sin embargo, la V de Cramér de \(0.183\)
indica que la magnitud de la asociación es pequeña.

El propósito permite reconocer segmentos con tasas diferentes, pero no
explica por sí solo la clasificación individual.

\subsubsection{¿La duración y el monto difieren entre las clases?}

Ambas variables presentan diferencias marginales.

Los malos créditos tienen una duración media aproximadamente
\(5.65\) meses mayor, con un tamaño de efecto:

\[
d=0.480.
\]

El monto también tiende a ser mayor, pero su tamaño de efecto es pequeño:

\[
r_{\mathrm{rb}}
=
0.110.
\]

La diferencia esencial aparece en el análisis multivariable: la duración
mantiene su asociación, mientras que el monto pierde significancia.

\subsubsection{¿Qué variables están asociadas con mayor riesgo?}

La evidencia conjunta identifica como señales principales:

\begin{itemize}
    \item mayor duración del crédito;

    \item mayor nivel de carga de la cuota;

    \item estado de cuenta menos favorable, especialmente la categoría
    de saldo negativo utilizada como referencia;

    \item menores niveles de ahorro;

    \item determinados perfiles del historial crediticio;

    \item menor estabilidad residencial frente a la vivienda propia.
\end{itemize}

La interpretación debe realizarse sobre el perfil completo. Ninguna
variable ofrece una separación perfecta entre buenos y malos créditos.

\subsubsection{¿Cómo cambia la clasificación al modificar el umbral?}

Con el umbral convencional:

\[
t=0.50,
\]

el modelo obtiene una exactitud de \(74.0\,\%\), pero identifica solo el
\(46.7\,\%\) de los malos créditos.

Bajo la relación de costos \(5:1\), el umbral teórico es:

\[
t^{*}
=
\frac{1}{6}
\approx
0.1667.
\]

Al aplicarlo, el recall aumenta a:

\[
88.9\,\%,
\]

y los falsos negativos disminuyen de \(48\) a \(10\). A cambio, los
falsos positivos aumentan y la exactitud disminuye al \(60.7\,\%\).

El costo total se reduce de:

\[
270
\quad\text{a}\quad
158,
\]

lo que representa una disminución aproximada del:

\[
41.5\,\%.
\]

La modificación del umbral no mejora la capacidad discriminativa del
modelo. Cambia la política de decisión y el equilibrio entre los dos
tipos de error.

\subsection{Coherencia entre los enfoques estadísticos}

Los resultados frecuentistas y bayesianos presentan una alta
consistencia.

Para la tasa global de mal crédito:

\[
\widehat{p}_{\mathrm{frecuentista}}
=
0.3000,
\]

mientras que, con el prior uniforme:

\[
\widehat{p}_{\mathrm{MAP}}
=
0.3000.
\]

El intervalo de Wilson y el intervalo creíble bajo
\(\operatorname{Beta}(1,1)\) coinciden numéricamente al nivel de redondeo:

\[
[0.2724,\;0.3291].
\]

Para la duración, el estimador frecuentista fue:

\[
\widehat{\beta}_{\mathrm{dur,MLE}}
=
0.03825,
\]

y la media posterior:

\[
\operatorname{E}
\left[
    \beta_{\mathrm{dur}}
    \mid
    \mathrm{datos}
\right]
=
0.04084.
\]

Las diferencias son pequeñas. Esta concordancia es esperable debido al
tamaño de la muestra y al uso de distribuciones previas débiles.

La similitud numérica no elimina la diferencia conceptual. La inferencia
frecuentista describe propiedades del procedimiento de estimación bajo
muestreo repetido. La inferencia bayesiana expresa incertidumbre
posterior condicionada al prior y a los datos observados.

\subsection{Predicción y utilidad operativa}

La regresión logística alcanzó:

\[
\operatorname{AUC}_{\mathrm{test}}
=
0.7740,
\]

y una media de validación cruzada de:

\[
\overline{\operatorname{AUC}}_{\mathrm{CV}}
=
0.7685
\pm
0.0221.
\]

La proximidad entre ambos resultados indica una capacidad discriminativa
estable dentro del esquema de evaluación utilizado.

Sin embargo, una AUC cercana a \(0.77\) no implica una clasificación
perfecta ni determina automáticamente una política crediticia adecuada.
El modelo ordena razonablemente los perfiles, pero existe un solapamiento
considerable entre las clases.

La utilidad operativa depende de tres niveles diferentes:

\begin{enumerate}
    \item \textbf{Estimación del riesgo:} producción de una probabilidad
    individual de mal crédito;

    \item \textbf{Discriminación:} capacidad de ordenar correctamente
    casos de mayor y menor riesgo;

    \item \textbf{Decisión:} elección de un umbral compatible con los
    costos y objetivos de la institución.
\end{enumerate}

Un modelo puede mostrar una discriminación estable y, al mismo tiempo,
producir decisiones costosas cuando se utiliza un umbral inadecuado.

\subsection{Implicancias para una política crediticia}

Los resultados sugieren que las operaciones de larga duración y elevada
carga de cuota requieren una revisión más cuidadosa. Esto no implica que
deban rechazarse automáticamente. Una política de gestión podría
considerar:

\begin{itemize}
    \item reducir el plazo cuando la capacidad de pago sea limitada;

    \item solicitar garantías adicionales en operaciones de mayor riesgo;

    \item ajustar el monto o la estructura de cuotas;

    \item establecer una revisión manual para probabilidades cercanas al
    umbral;

    \item incorporar información financiera más reciente antes de la
    decisión final.
\end{itemize}

El monto no debería utilizarse aisladamente como argumento de riesgo. Su
asociación marginal desaparece después de considerar la duración y otras
características.

Del mismo modo, categorías personales o históricas no deben convertirse
automáticamente en reglas de exclusión. Su utilización debe someterse a
criterios de pertinencia, equidad y regulación.

\subsection{Respuesta integrada final}

La probabilidad de mal riesgo crediticio se asocia principalmente con la
combinación de plazos prolongados, mayor carga de la cuota y señales menos
favorables de liquidez o estabilidad financiera.

La duración es el factor numérico más consistente: aparece en el análisis
exploratorio, la comparación frecuentista, la regresión logística y la
posterior bayesiana. El estado de cuenta, los ahorros, el historial
crediticio y la vivienda complementan esta señal.

El monto y la edad muestran diferencias marginales, pero no conservan
evidencia condicional dentro del modelo final. El propósito está asociado
globalmente con el riesgo, aunque su magnitud es pequeña y su utilidad
principal es la segmentación.

Finalmente, el desempeño del sistema no depende únicamente de la calidad
del modelo. La selección del umbral modifica sustancialmente los errores
y los costos. Bajo la matriz \(5:1\), el umbral \(1/6\) resulta más
coherente que \(0.50\), porque prioriza la detección de malos créditos y
reduce el costo total observado.

En consecuencia, el modelo debe entenderse como una herramienta de apoyo
a la decisión: combina información sobre el perfil y la operación,
produce una estimación de riesgo y permite adaptar la clasificación a los
costos del problema. No constituye una explicación causal ni un sistema
listo para despliegue sin validación externa, recalibración y revisión de
equidad.

% Contenido previsto:
% - Respuesta conjunta a las cinco preguntas del Grupo 5.
% - Diferencia entre asociación marginal y condicional.
% - Convergencia y divergencia entre resultados frecuentistas y bayesianos.
% - Implicancias operativas del umbral y los costos.
